\chapter{Competitor Rules}

\section{General Road Racing Rules}

\subsection{Safety}
Riders must wear shoes, gloves and a helmet (see definitions in chapter \ref{chap:general_definitions}).
Knee pads and elbow pads are suggested optional safety gear.

Personal music systems are not allowed for any races on public roads where there may be motorized traffic.

Water and food are the responsibility of the rider.
It is recommended riders carry their own water on races longer than 10k.
Hosts may offer food and water stations at their discretion.

\subsection{Unicycles}

\begin{enumerate}[leftmargin=1.5em,labelsep=*]
	\item Riders may use different unicycles for different road events, as long as they comply with rules for the events in which they are entered.
	\item Unicycle racing is divided into two categories: standard and unlimited.
	In the standard category, unicycles have a standardised wheel size and fixed transmission.
	The recognised classes in the standard category are the 16 class, 20 class, 24 class and 29 class.
	In the unlimited category unicycles are not-standardised by wheel size.
	The recognised class in the unlimited category is the unlimited class.
	\item There is an allowable tire diameter range, minimum crank arm length and transmission for each unicycle class:
	\begin{longtable}{|p{2.8cm}|p{4.6cm}|p{3.6cm}|p{2.5cm}|}
		\hline
		\textbf{Unicycle Class} & \textbf{Diameter Range} & \textbf{Min Crank Length} & \textbf{Transmission}\\
		\hline
		16 Class & 0 -- 418mm & 89mm & regular \\
		\hline
		20 Class & more than 418mm -- 518mm & 100mm & regular \\
		\hline
		24 Class & more than 518mm -- 618mm & 125mm & regular \\
		\hline
		29 Class & more than 618mm -- 778mm & No limit & regular \\
		\hline
		Unlimited Class & No limit & No limit & unlimited \\
		\hline
	\end{longtable}
	\item For Unicon and continental championships, both standard and unlimited competition categories are mandatory for time-trials, criterium and fixed/free-distance races. 
    For smaller regional and national championship events, the two competition categories are strongly recommended.  Exceptions apply (see section 3B.2.4.1 )
	\item For fixed/free distance races:
	\begin{enumerate}[label*=\arabic*,leftmargin=1.5em,labelsep=*]
		\item If the distance is greater than 10km, the standard category is the 29 Class, and the Unlimited category is the Unlimited Class.
		\item If the distance is 10km or less, the standard category is the 24 Class and the unlimited category is the unlimited class.
		For age group competition, there should be a 20 Class for riders under 11 years old.
		The youngest age group for 24 Class wheels should have a minimum age of 0, so riders 10 years and younger also have the option of racing in the 24 Class 
        (e.g. 0-8 on 20 Class, 9-10 on 20 Class, 0-13 on 24 Class).
	\end{enumerate}
	\item For time-trials, the standard category is the 29 class and for the unlimited category is the unlimited class. 
    For criterium, the standard category is the 24 class and the unlimited category the unlimited class.
	\item A unicycle can be ridden if it is smaller than the class in which it is entered
	(e.g.\ 20 Class unicycle is allowed in a 24 Class race).
	However, if a race is over 20km, 24 Class and smaller unicycles are not permitted unless exempted by the race director.
	\item Unicycles must be checked for compliance within their wheel class (wheel diameter, crank length and gearing), with the tire pressure that will be used in the race.
	Preferably, this check is carried out immediately before the race.
	Crank arm length is measured from the center of the wheel axle to the center of the pedal axle.
	Longer sizes may be used.
\end{enumerate}

\subsection{Rider Identification}
%TODO Nummerierung der einzelnen Absätze
\subsubsection{Race Numbers}
	Riders must wear their race number so that it is clearly visible during the race and when crossing the finish line.
	The recommended position for the start number is on the back.
	Organizers may provide special race numbers for the helmet, arms or hips, so that the riders can be identified from the side during the race.
	
	Riders must use the officially provided race number unmodified in any manner.
	Numbers should not be folded, trimmed, or otherwise defaced.
	Referee approval must be sought to modify a race number if it cannot otherwise be attached securely due to hydration pack, rider physique or posture when riding.
	Lost or damaged race numbers must be replaced with approval by referee.

\subsubsection{Chips for Electronic Timing}
	Riders may be required to wear an electronic timing chip, as specified by the timing personnel.
	This chip may be attached, among other methods, as a disposable tag to the race number, which may have to be worn on the chest.

\subsection{Riders Must Be Ready}

Riders must be ready when called for their races.
Riders not at the start line when their race begins may lose their chance to participate.
The Starter will decide when to stop waiting, remembering to consider language barriers, and the fact that some riders may be slow because they are helping run the convention.

\subsection{Protests}
%TODO Abgleich mit Kapitel 1 und ggf. kürzen
	\begin{enumerate}[leftmargin=1.5em,labelsep=*]
		\item Protests must be filed on an official form.
		Mistakes in paperwork, inaccuracies in placing, and interference from other riders or other sources are all grounds for protests. 
        All Referee decisions are final, and cannot be protested.
		\item The default protest time (counting from the posting of results) depends on the type of event and the time after the competition within the results are published.
		The default protest time can be extended or shortend up to the minimum by the Race Director.
		Every deviation from the default protest time has to be clearly announced when the results are posted, including stating the protest deadline on the results list itself.
		The protest time may be extended for riders who have to be in other races during the protest period.
		\begin{enumerate}[label*=\arabic*,leftmargin=1.5em,labelsep=*]
			\item For a large event such as Unicon or continental championships, for all results published within 90 minutes 
            after the end of the competiotion, the default protest time is 120 minutes, the minimum is 30 minutes.
			For results published more than 90 minutes after the end of the competition but before 8 p.m. of the same day, the minimum protest time is 120 minutes.
			For all results published after 8 p.m. on the same day or on another day, the protest deadline is 12 p.m. (noon time) on the day after publication.
			\item For smaller events, for all results published within 90 minutes after the end of the competition, the default protest time is 60 minutes, the minimum is 15 minutes.
			For results published more than 90 minutes after the end of the competition, the minimum protest time is 60 minutes.
		\end{enumerate}
		\item All protests will be acknowledged within 30 minutes from the time they are received, and an effort will be made to settle the issue within those 30 minutes.
	\end{enumerate}

\section{Road Racing Disciplines}

\subsection{Fixed Distance Races}
	\begin{enumerate}[leftmargin=1.5em,labelsep=*]
		\item Fixed Distance Races are drafting races, with a wave or mass start, where riders races against each other over a defined and precisely measured distance.
		\item The officiall recognized distances are
		\begin{enumerate}[label=\roman*,leftmargin=2.25em,labelsep=*]
			\item 10 km
			\item Marathon (42.195 km)
		\end{enumerate}
	\end{enumerate}

\subsection{Free Distance Races}
	\begin{enumerate}[leftmargin=1.5em,labelsep=*]
		\item Free Distance Races are drafting races, with a wave or mass start, where riders race against each other in a non-recognized fixed distance.
		\item A Free Distance Race can be any race with a distance that is greater or lesser than 3 \% from an officially recognized Fixed Distance.
		For Free Distance Races longer than a Marathon, a mass start is preferred.
		\item The approximate distance of the race must be announced at the time of registration, to help riders plan and prepare accordingly.
		\item There is no requirement for the course of a Free Distance Race to be officially measured.
		However, once the final course has been confirmed, the distance must be communicated to riders with at least kilometer-level precision.
		Organizers are encouraged to share a GPS track of the confirmed course whenever possible.
	\end{enumerate}

\subsection{Criterium}
	\begin{enumerate}[leftmargin=1.5em,labelsep=*]
		\item A Criterium is a short drafting race with distances of about 5 km to 10 km.
		It can be held around city block(s) or within a large parking lot.
		The courses should have left and right turns and multiple laps with a recommended lap length of 500 m to 1000 m.
		\item A Criterium can be run in one of the following two variants:
		\begin{enumerate}[label*=\arabic*,leftmargin=1.5em,labelsep=*]
			\item As a set distance criterium, where all riders will complete the same number of laps.
			\item As a time-based criterium, in which all riders ride one more lap after a set time and then everyone finishes the race.
		\end{enumerate}
	\end{enumerate}
	
\subsubsection{Set Distance Criterium}
	\begin{enumerate}[leftmargin=1.5em,labelsep=*]
		\item The number of laps should be three to ten laps and must be announced clearly to riders before the start of the race.
		\item Each rider is responsible for counting their laps; organizers are not responsible for riders who do not finish the race due to not completing the required number of laps.
	\end{enumerate}

\subsubsection{Time-Based Criterium}
	\begin{enumerate}[leftmargin=1.5em,labelsep=*]
		\item Using the time from the top rider's first two laps, the referee will determine how many laps could be completed in the desired time limit.
		From this point on, the number of remaining laps (for the leaders) will be displayed and this will be used to determine when finish of the race occurs.
		A bell will be rung with one lap to go.
		\item Lapped riders in the race will all finish on the same lap as the leader and will be placed according to the number of 
        laps they are down and then their position at the finish.
	\end{enumerate}

\subsection{Individual Time Trial}
	\begin{enumerate}[leftmargin=1.5em,labelsep=*]
		\item An Individual Time Trial is a non-drafting race, with an individual start for each rider, where each rider races against the clock.
		\item An Individual Time Trial can be run in one of the following two variants, either as part of an 
        official unicycling competition or as independent events, such as record attempts:
		\begin{enumerate}[label*=\arabic*,leftmargin=1.5em,labelsep=*]
			\item As a distance-based time trial, in which a given distance is completed as fast as possible.
			\item As a time-based time trial, in which the greatest possible distance is completed in a fixed amount of time.
		\end{enumerate}
	\end{enumerate}

\subsubsection{Distance-based Time Trials}
	\begin{enumerate}[leftmargin=1.5em,labelsep=*]
		\item These events may be held either during an official competition or independently, including for official record purposes.
		\begin{enumerate}[label*=\arabic*,leftmargin=1.5em,labelsep=*]
			\item In official competitions, there is no fixed distance requirement for time trial course.
			Organizers may choose any course length that fits the format of their event.
			It is recommended to use a point-to-point course without loops, in order to minimize interactions between participants.
			The course should allow riders to overtake safely with sufficient space over most of the course.
			\begin{enumerate}[label*=.\arabic*,leftmargin=2.25em,labelsep=*]
				\item Riders must start one by one at fixed intervals, typically every 20 to 30 seconds.
				The starting order should be based on estimated performance, with slower riders starting first and faster riders starting later.
%TODO For consistency of English language text, please change 'Brutto' to 'Gross' or 'Gun' time and 'Netto' to 'Net' time.
				\item The organizer may choose between brutto timing (elapsed time from the official start time of each rider) 
                or netto timing (measured using a chip or electronic device triggered upon crossing the start line).
				\item If a 10 km fixed-distance course is used for a 24" Standard race, it may also be used for a time trial.
				Two options are possible:
				\begin{enumerate}[label=\roman*,leftmargin=2.25em,labelsep=*]
					\item Preferred option: Hold the 24" Standard fixed-distance race on a different day.
					This allows the time trial to be held with both 29" Standard and Unlimited categories.
					\item Alternative option: If it is not possible to separate the events across different days, 
                    the time trial may include only an Unlimited category and should take place just before or 
                    after the 24" Standard race on the same course.
				\end{enumerate}
			\end{enumerate}
			\item For time trials held outside of official competitions, such as those organized for the purpose of record attempts, 
            it is recommended to use a closed course on a road or track.
			The course must be measured and verified in accordance with the WR guidelines.
		\end{enumerate}
		\item Official formats recognized for the purpose of establishing official records are
		\begin{enumerate}[label=\roman*,leftmargin=2.25em,labelsep=*]
			\item 10 km time trials and
			\item 100 km time trials and
			\item 100 miles time trials.
		\end{enumerate}
	\end{enumerate}

\subsubsection{Time-based Time Trials}
	\begin{enumerate}[leftmargin=1.5em,labelsep=*]
		\item Time-based Time Trials must take place on a closed circuit, either on a road or a track.
		The course must be designed so that the start and finish occur at the exact same location.
		These events are not held as part of official competitions, but are recognized for the purpose of establishing official records.
		\item Official formats recognized for the purpose of establishing official records are
		\begin{enumerate}[label=\roman*,leftmargin=2.25em,labelsep=*]
			\item 1-hour time trials and
			\item 24-hour time trials.
		\end{enumerate}
	\end{enumerate}

\subsection{Hill Climb}
	\begin{enumerate}[leftmargin=1.5em,labelsep=*]
		\item A Hill Climb is a unicycle road race held on a sustained uphill course, 
        designed to test the riders' climbing ability rather than their top speed on flat or descending terrain.
		\item A Hill Climb must take place in a region with appropriate topography, featuring a continuous ascent over several kilometers.
		Flat or downhill sections should be avoided, especially at the start and the end of the course.
		A flat section may be used if required for a safe and practical race start.
		The course must either have an average gradient of at least 6 \%, or a total elevation gain of at least 850 meters.
		\item There is no fixed maximum distance for a Climbing Road Race.
		The course length should be determined based on local topography, aiming for meaningful and challenging durations.
		Organizers are encouraged to consider a minimum distance of 5 km.
		\item If topography allows, organizers may offer two course options: a longer "Expert" course and a shorter "Beginner" course.
		Both courses should follow the same climb whenever possible.
		\item The majority of the course must be on paved road surfaces.
		Short unpaved sections are allowed only if they are smooth and rideable.
		The entire course must be rideable by elite riders without requiring dismounts due to excessive steepness or technical difficulty.
		\item Intentionally running to gain time is not permitted.
		\item A Hill Climb is a single unlimited-category competition, meaning that all riders, including standard category, 
        start together in the same race without separation by unicycle category.
		When the event is held at an UNICON, additional awards shall be given to the top three males and top three females using 29" standard unicycles.
	\end{enumerate}


\section{Racing Rules}

\subsection{The Start}

\begin{enumerate}[leftmargin=1.5em,labelsep=*]
 	\item Riders start with the fronts of their tires (forward most part of wheel) behind the edge of the starting line that is farthest from the finish line.
	 \begin{enumerate}[label*=\arabic*,leftmargin=1.5em,labelsep=*]
	 	\item With an individual start, the rider may start mounted, holding onto a starting post or other support.
	 	The rider may place the starting post or other support in the location most comfortable for them.
	 	The rider may mount after the start signal, if they wish.
	 	The rider has to start directly behind the Start line.
 		\item With an heat or mass start, riders may start mounted, holding onto a starting post or other support, or onto each other.
 		Riders may mount after the start signal, if they wish.
	 \end{enumerate}
	 \item Rolling starts are not permitted in any race.
	 Riders may lean before the start, but their wheels may not move forward during the start beeps or counting down.
	 Rolling back should be avoided.
	 \item All commands of the starter are to be given in English at Unicon or international competitions. At other competitions, English is optional.
	 \item For heat or mass starts, the starter must announce the last remaining minute before the start.
	 About 10 seconds before the start, the starter should give the command “Ready”.
	 After the command “Ready”, all riders must move to their starting position.
	 As soon as the Starter is satisfied that all riders are steady in the correct starting position, they gives the command “Attention” and starts the race.
	 This can be done by a start sequence as follows:
	 \begin{enumerate}[label*=\arabic*,leftmargin=1.5em,labelsep=*]
	 	\item Usually, a start-beep apparatus is used.
	 	This provides a six-count start: ``beep - beep -beep - beep - beep - buup!''.
	 	The timing between (the start of) successive beeps is one second.
	 	The first five beeps have all the same sound frequency.
	 	The final tone (buup) has a higher frequency, so that the competitors can easily distinguish this tone from the rest.
	 	The proper moment to start is the beginning of the final tone.
	 	\item As an alternative, the Starter will give a three-count start before firing a starting gun on the fourth count.
	 	Example: ``One, two, three, BANG!''.
	 	The time between (the start of) each of these elements must be the same, and should approximately 1 second.
	 \end{enumerate}
	 Both variants allow the rider to start leaning ahead of the “buup/BANG”, for an exact and predictable start.
	 It is recommended to use one or the other of those two options for all races in a competition if possible.
	 The option to be used must be announced in advance of the competition.
	 \item If the Starter is not satisfied that all is ready for the start to proceed after they gave the command ``Ready'' and the riders are on their starting position or they otherwise abort the start, the command should be ``Go Back''.
	 If a start-beep apparatus is used and the start sequence is already started the start should be aborted by blowing a whistle or other clear and predefined signal.
	 Where a rider in the judgement of the Starter, after the command ``Ready'',
	 \begin{enumerate}[label*=\arabic*,leftmargin=1.5em,labelsep=*]
	 	\item causes the start to be aborted, for instance by dismounting, without a valid reason (such reason to be evaluated by the Starter); or
	 	\item does not place themselves in their final starting position at once and without delay; or
	 	\item disturbs other riders in the race through sound, movement or otherwise,
	 \end{enumerate}
	 the Starter must abort the start.
	 The Starter may warn the rider for improper conduct (disqualify in case of repeated infringement of the Rule).
	 However, when an extraneous reason was considered to be the cause for aborting the start, or the Referee does not agree with the Starters decision, 
     no rider gets warned or disqualified.
	 This decision must be clearly indicated to the riders.
\end{enumerate}

\subsubsection{False Start}

\begin{enumerate}[leftmargin=1.5em,labelsep=*]
	\item A false start occurs if a rider's wheel moves forward before the start signal.
%TODO Abgleich mit Track Regel - "any"
	\item If a heat has to be restarted, the Starter will immediately recall the riders, for example by blowing a whistle or any other clear and predefined signal.
	Any warning or disqualification resulting from this must be clearly indicated to the riders in question.
	\item There are three options on how to deal with false starts:
	\begin{enumerate}[label*=\arabic*,leftmargin=1.5em,labelsep=*]
		\item \textbf{One False Start Allowed Per Heat:}
		The use of this option is strongly discouraged when no electronic false start monitoring system is used.
		In case of a false start, the heat is restarted.
		After the first false start of a particular heat, all riders receive a warning and may start again.
		Thereafter, any rider(s) causing a false start are disqualified for this event.
		Only the earliest false starting rider gets assigned this false start and the associated disqualification.
		\item \textbf{One False Start Allowed Per Rider:}
		In case of a false start, the heat is restarted.
		After the first false start of a particular rider in a heat, the rider in question receives a warning and may start again.
		Any rider(s) causing their personal second false start are disqualified for this event.
		Only the earliest false starting rider gets assigned this
        false start and the associated warning or disqualification.
		\item \textbf{Time Penalty:}
		In case of a false start, the heat is not restarted.
		If a false start occurs by one or multiple riders, these riders receive a time penalty (such as 10 seconds).
	\end{enumerate}
	It is recommended to use one or the other of those three options for all races in a competition if possible.
	The option to be used must be announced in advance of the competition.
\end{enumerate}

\subsection{The Race}

\subsubsection{Illegal Riding}

Illegal riding includes intentionally interfering in any way with another rider, deliberately crossing in front of another rider to prevent him or her from moving on, deliberately blocking another rider from passing, or distracting another rider with the intention of causing a dismount.

\subsubsection{Passing}

An overtaking rider must pass on the outside, unless there is enough room to safely pass on the inside.
Riders passing on the inside are responsible for any fouls that may take place as a result.
No physical contact between riders is allowed.
The slower rider must maintain a reasonably straight course, and not interfere with the faster rider.

\subsubsection{Dismounts}

Dismounting and remounting is allowed.
If a rider is forced to dismount due to a fall by the rider immediately in front, it is considered part of the race, and both riders should remount and continue.

\subsubsection{Repair, Change, or Replace a (Broken) Unicycle \label{subsec:road_unicycle-changes}}

	\begin{enumerate}[leftmargin=1.5em,labelsep=*]
		\item Riders are allowed to modify or adjust their unicycle during the race.
		Assistance from others is permitted at any time, including tools and hands-on help.
		\item Replacement of the unicycle or wheel is permitted only in the event of mechanical problems or damage.
		The replacement unicycle must belong to the same category as the original and should ideally have same characteristics (wheel size, crank length, hub).
		The Referee must confirm that the situation was unplanned and was indeed accidental.
		\item The rider may continue the course on foot (walking, not running) with the broken unicycle.
		If they leave the course, they must reenter the course at or before the point where they exited from the course.
		While off-course, they may run or use any form of transportation.
	\end{enumerate}

\subsubsection{Drafting}
	\begin{enumerate}[leftmargin=1.5em,labelsep=*]
		\item Drafting is riding close behind another rider or vehicle to reduce air resistance and benefit from the slipstream.
		It is considered drafting when the distance is less than 5 meters front/back or 1.5 meters laterally.
		\item In all road races drafting behind non-unicyclist participants, non-participants and vehicles is prohibited at all times.
		Additinally in all
		\begin{enumerate}[label*=\arabic*,leftmargin=1.5em,labelsep=*]
			\item non-drafting races, drafting behind other unicyclist participants is forbidden at any time.
			\item drafting races, drafting between unicyclist participants in the same race is permitted unless otherwise specified in the event rules.
		\end{enumerate}
		\item Race officials may evaluate drafting infractions based on rider behavior, timing, and positioning.
	\end{enumerate}

\subsection{The Finish}

Finish times are determined when the front of the tire first crosses the vertical plane of the nearest edge of the finish line.

Riders are always timed by their wheels, not by outstretched bodies.
If riders do not cross the line in control, they are awarded a 5 second penalty to their time.
``Control'' is defined by the front of the wheel crossing the vertical finish plane (as defined above) with the rider having both feet on the pedals.
(Note: a rider is not considered in control if the unicycle crosses the finish line independent of the rider.
The finish time is still measured by when the wheel crosses the vertical finish plane and the 5 second penalty is applied.)

In the case where a rider is finishing with a broken unicycle, the rider must bring at minimum the wheel to the finish line, and time is still taken when the wheel crosses the finish line.
The 5 second penalty is applied.
